\documentclass[12pt,a4paper]{article}

\usepackage[T2A]{fontenc}
\usepackage[utf8]{inputenc}
\usepackage[russian]{babel}
\usepackage{amsmath, amssymb, amsthm, mathtools}
\usepackage{array}
\usepackage{booktabs}
\usepackage{enumitem}
\usepackage{tikz}
\usetikzlibrary{positioning}
\usepackage{listings}
\usepackage{xcolor}
\usepackage{geometry}

\geometry{left=25mm,right=20mm,top=20mm,bottom=20mm}

\newtheorem{theorem}{Теорема}
\newtheorem{definition}{Определение}
\newtheorem{example}{Пример}

\lstset{
    basicstyle=\ttfamily\small,
    keywordstyle=\color{blue},
    commentstyle=\color{gray},
    frame=single,
    breaklines=true,
    columns=fullflexible
}

\title{Билет №70 \\ \small Операционные средства управления сетями. Эталонная модель взаимодействия открытых систем ISO/OSI. Маршрутизация и управление потоками данных в сети. (ИТМО)}
\author{Сериков Владислав Александрович}
\date{16 мая 2026}

\begin{document}

\maketitle
\tableofcontents
\newpage

\section{Компьютерная сеть и задачи сетевого управления}

\begin{definition}
\textbf{Компьютерная сеть} --- это совокупность узлов, каналов связи, сетевого оборудования и программных средств, обеспечивающих обмен данными между вычислительными системами.
\end{definition}

В общем случае сеть включает:
\begin{itemize}
    \item \textbf{конечные узлы}: рабочие станции, серверы, мобильные устройства, IoT-устройства;
    \item \textbf{промежуточные узлы}: коммутаторы, маршрутизаторы, межсетевые экраны, точки доступа;
    \item \textbf{каналы связи}: медные, оптические, радиоканалы;
    \item \textbf{сетевые протоколы}: Ethernet, IP, TCP, UDP, OSPF, BGP, SNMP и др.;
    \item \textbf{средства управления}: операционные системы, сетевые службы, системы мониторинга, средства конфигурационного управления.
\end{itemize}

\begin{definition}
\textbf{Управление сетью} --- это совокупность процессов, протоколов и программных средств, предназначенных для наблюдения за состоянием сети, настройки сетевых устройств, обнаружения неисправностей, анализа производительности, обеспечения безопасности и планирования развития сети.
\end{definition}

Классическая модель функциональных областей управления сетью часто описывается аббревиатурой \textbf{FCAPS}:

\begin{center}
\begin{tabular}{|p{0.22\textwidth}|p{0.65\textwidth}|}
\hline
\textbf{Область} & \textbf{Содержание} \\
\hline
Fault Management & Обнаружение, локализация и устранение отказов. \\
\hline
Configuration Management & Настройка оборудования, ведение конфигураций, инвентаризация. \\
\hline
Accounting Management & Учёт использования ресурсов, тарификация, квоты. \\
\hline
Performance Management & Измерение загрузки каналов, задержек, потерь, пропускной способности. \\
\hline
Security Management & Контроль доступа, аудит, шифрование, управление политиками безопасности. \\
\hline
\end{tabular}
\end{center}

\section{Операционные средства управления сетями}

\subsection{Сетевые функции операционных систем}

Современная операционная система выполняет не только локальное управление процессами, памятью и файлами, но и содержит сетевую подсистему.

Основные сетевые функции ОС:
\begin{enumerate}
    \item реализация стека протоколов TCP/IP;
    \item управление сетевыми интерфейсами;
    \item маршрутизация пакетов;
    \item фильтрация и межсетевое экранирование;
    \item поддержка сетевых служб: DNS, DHCP, SSH, HTTP, NTP и др.;
    \item сбор статистики и журналирование сетевых событий;
    \item поддержка виртуальных сетей и туннелей;
    \item управление качеством обслуживания и очередями пакетов.
\end{enumerate}

\subsection{Сетевой стек ОС}

Сетевой стек операционной системы обычно включает несколько уровней:

\begin{center}
\begin{tabular}{|c|p{0.6\textwidth}|}
\hline
\textbf{Уровень} & \textbf{Примеры функций} \\
\hline
Прикладной & Сокеты, системные вызовы, сетевые службы. \\
\hline
Транспортный & TCP, UDP, управление соединениями и портами. \\
\hline
Сетевой & IP-адресация, маршрутизация, ICMP. \\
\hline
Канальный & Ethernet, ARP, драйверы сетевых карт. \\
\hline
Физический & Передача битов по среде, реализуется аппаратурой. \\
\hline
\end{tabular}
\end{center}

С точки зрения ОС, сетевой обмен часто реализуется через абстракцию \textbf{сокета}.

\begin{definition}
\textbf{Сокет} --- программный интерфейс, представляющий конечную точку сетевого взаимодействия. Обычно сокет задаётся адресом узла, номером порта и транспортным протоколом.
\end{definition}

Например, TCP-соединение можно описать пятёркой:
\[
(\text{IP}_{src}, \text{port}_{src}, \text{IP}_{dst}, \text{port}_{dst}, \text{protocol}).
\]

\subsection{Средства настройки сети в ОС}

Типичные операции управления сетью:
\begin{itemize}
    \item назначение IP-адреса и маски подсети;
    \item настройка шлюза по умолчанию;
    \item настройка DNS-серверов;
    \item просмотр таблицы маршрутизации;
    \item добавление статических маршрутов;
    \item включение или отключение интерфейсов;
    \item настройка правил межсетевого экрана;
    \item мониторинг соединений и трафика.
\end{itemize}

Примеры команд в UNIX-подобных ОС:

\begin{lstlisting}
ip addr show
ip route show
ip route add 10.10.0.0/16 via 192.168.1.1
ss -tulpen
ping 8.8.8.8
traceroute example.org
tcpdump -i eth0
\end{lstlisting}

Примеры команд в Windows:

\begin{lstlisting}
ipconfig /all
route print
netstat -ano
ping 8.8.8.8
tracert example.org
nslookup example.org
\end{lstlisting}

\subsection{Мониторинг и диагностика}

Основные показатели состояния сети:
\begin{itemize}
    \item доступность узлов;
    \item задержка передачи;
    \item вариация задержки, или jitter;
    \item процент потерь пакетов;
    \item загрузка каналов;
    \item число ошибок на интерфейсах;
    \item число TCP-повторных передач;
    \item состояние таблиц маршрутизации;
    \item загрузка CPU и памяти сетевых устройств.
\end{itemize}

Типичные диагностические средства:
\begin{itemize}
    \item \texttt{ping} --- проверка доступности узла с помощью ICMP Echo;
    \item \texttt{traceroute} / \texttt{tracert} --- определение маршрута до узла;
    \item \texttt{netstat}, \texttt{ss} --- просмотр сетевых соединений;
    \item \texttt{tcpdump}, Wireshark --- анализ сетевого трафика;
    \item SNMP-системы мониторинга --- сбор информации с сетевых устройств;
    \item syslog --- централизованный сбор журналов;
    \item NetFlow/IPFIX --- учёт потоков трафика.
\end{itemize}

\section{Протоколы и архитектуры управления сетью}

\subsection{SNMP}

\begin{definition}
\textbf{SNMP} --- Simple Network Management Protocol, протокол прикладного уровня, предназначенный для мониторинга и управления сетевыми устройствами.
\end{definition}

Архитектура SNMP включает:
\begin{itemize}
    \item \textbf{manager} --- управляющая станция;
    \item \textbf{agent} --- программа на управляемом устройстве;
    \item \textbf{MIB} --- Management Information Base, база управляемых объектов;
    \item \textbf{OID} --- Object Identifier, уникальный идентификатор объекта в MIB.
\end{itemize}

Основные операции SNMP:
\begin{itemize}
    \item \texttt{GET} --- получить значение объекта;
    \item \texttt{GETNEXT} --- получить следующий объект в дереве MIB;
    \item \texttt{GETBULK} --- получить сразу множество объектов;
    \item \texttt{SET} --- изменить значение объекта;
    \item \texttt{TRAP} --- асинхронное уведомление от агента менеджеру;
    \item \texttt{INFORM} --- подтверждаемое уведомление.
\end{itemize}

Пример:
\[
\text{manager} \xrightarrow{\text{GET ifInOctets.2}} \text{agent}
\]
\[
\text{agent} \xrightarrow{\text{value = 12500000}} \text{manager}.
\]

\subsection{Syslog}

\begin{definition}
\textbf{Syslog} --- механизм передачи и централизованного хранения сообщений журналов от сетевых устройств и серверов.
\end{definition}

Сообщение syslog обычно содержит:
\begin{itemize}
    \item временную метку;
    \item имя устройства;
    \item компонент или процесс;
    \item уровень важности;
    \item текст сообщения.
\end{itemize}

Уровни важности:
\[
\text{Emergency} > \text{Alert} > \text{Critical} > \text{Error} > \text{Warning} > \text{Notice} > \text{Info} > \text{Debug}.
\]

\subsection{NetFlow и IPFIX}

\begin{definition}
\textbf{Поток трафика} --- совокупность пакетов, имеющих общий набор признаков, например одинаковые адреса источника и назначения, порты и протокол.
\end{definition}

Обычно поток задаётся пятёркой:
\[
(srcIP, dstIP, srcPort, dstPort, protocol).
\]

NetFlow/IPFIX позволяют оценивать:
\begin{itemize}
    \item какие узлы создают наибольшую нагрузку;
    \item какие приложения используют канал;
    \item какие направления трафика преобладают;
    \item есть ли аномалии или признаки атак.
\end{itemize}

\subsection{SDN и централизованное управление}

\begin{definition}
\textbf{SDN} --- Software Defined Networking, подход к построению сетей, при котором плоскость управления отделяется от плоскости передачи данных.
\end{definition}

В традиционной сети каждое устройство самостоятельно принимает решения о маршрутизации и коммутации. В SDN логика управления выносится в контроллер.

\[
\text{Приложения управления}
\longleftrightarrow
\text{SDN-контроллер}
\longleftrightarrow
\text{сетевые устройства}
\]

Преимущества SDN:
\begin{itemize}
    \item централизованная политика управления;
    \item программируемость сети;
    \item автоматизация конфигураций;
    \item гибкое управление потоками;
    \item удобство построения виртуальных сетей.
\end{itemize}

\section{Эталонная модель взаимодействия открытых систем ISO/OSI}

\subsection{Назначение модели}

\begin{definition}
\textbf{Модель ISO/OSI} --- эталонная семиуровневая модель взаимодействия открытых систем, предназначенная для описания функций сетевого взаимодействия и стандартизации протоколов.
\end{definition}

Модель OSI не является конкретным стеком протоколов. Это концептуальная модель, позволяющая:
\begin{itemize}
    \item разделить сложную задачу сетевого взаимодействия на уровни;
    \item определить функции каждого уровня;
    \item описать взаимодействие соседних уровней;
    \item обеспечить независимость реализации уровней;
    \item упростить стандартизацию и диагностику.
\end{itemize}

\subsection{Основные понятия модели OSI}

\begin{definition}
\textbf{Уровень} --- логическая часть сетевой архитектуры, выполняющая определённый набор функций и предоставляющая сервис вышележащему уровню.
\end{definition}

\begin{definition}
\textbf{Протокол} --- набор правил взаимодействия одноимённых уровней на разных узлах.
\end{definition}

\begin{definition}
\textbf{Сервис} --- набор функций, предоставляемых уровнем вышележащему уровню.
\end{definition}

\begin{definition}
\textbf{Интерфейс} --- правила взаимодействия соседних уровней внутри одной системы.
\end{definition}

Важно различать:
\begin{itemize}
    \item \textbf{вертикальное взаимодействие} --- между соседними уровнями на одном узле;
    \item \textbf{горизонтальное взаимодействие} --- между одноимёнными уровнями разных узлов по протоколу.
\end{itemize}

\subsection{Инкапсуляция}

При передаче данных каждый уровень добавляет к данным служебную информацию.

\begin{definition}
\textbf{Инкапсуляция} --- процесс добавления заголовков и, возможно, трейлеров к данным при их прохождении сверху вниз по стеку протоколов.
\end{definition}

Упрощённо:
\[
\text{данные}
\rightarrow
\text{сегмент}
\rightarrow
\text{пакет}
\rightarrow
\text{кадр}
\rightarrow
\text{биты}.
\]

\begin{center}
\begin{tikzpicture}[node distance=0cm]
\node[draw, minimum width=7cm, minimum height=0.8cm] (app) {Данные приложения};
\node[draw, below=of app, minimum width=8cm, minimum height=0.8cm] (tcp) {TCP-заголовок + данные};
\node[draw, below=of tcp, minimum width=9cm, minimum height=0.8cm] (ip) {IP-заголовок + TCP-сегмент};
\node[draw, below=of ip, minimum width=10cm, minimum height=0.8cm] (eth) {Ethernet-заголовок + IP-пакет + FCS};
\end{tikzpicture}
\end{center}

\subsection{Семь уровней OSI}

\begin{center}
\begin{tabular}{|c|p{0.22\textwidth}|p{0.48\textwidth}|p{0.16\textwidth}|}
\hline
\textbf{№} & \textbf{Уровень} & \textbf{Функции} & \textbf{Примеры} \\
\hline
7 & Прикладной & Сетевые службы для приложений & HTTP, DNS, SMTP \\
\hline
6 & Представления & Форматы данных, кодирование, шифрование, сжатие & TLS, ASN.1, JPEG \\
\hline
5 & Сеансовый & Управление сеансами, диалогом, синхронизацией & RPC, NetBIOS \\
\hline
4 & Транспортный & Доставка между процессами, порты, надёжность, управление потоком & TCP, UDP \\
\hline
3 & Сетевой & Логическая адресация, маршрутизация, фрагментация & IP, ICMP, OSPF \\
\hline
2 & Канальный & Доставка в пределах канала, MAC-адреса, кадры, контроль ошибок & Ethernet, Wi-Fi, PPP \\
\hline
1 & Физический & Передача битов по среде & UTP, оптика, радио \\
\hline
\end{tabular}
\end{center}

\subsection{Физический уровень}

Физический уровень определяет:
\begin{itemize}
    \item электрические, оптические или радиосигналы;
    \item разъёмы и кабели;
    \item кодирование битов;
    \item частоты и уровни сигналов;
    \item скорость передачи.
\end{itemize}

Единица данных:
\[
\text{бит}.
\]

\subsection{Канальный уровень}

Канальный уровень обеспечивает передачу кадров между соседними узлами одного канала или сегмента сети.

Основные функции:
\begin{itemize}
    \item формирование кадров;
    \item физическая адресация, например MAC-адреса;
    \item обнаружение ошибок;
    \item управление доступом к среде;
    \item иногда --- исправление ошибок.
\end{itemize}

Единица данных:
\[
\text{кадр}.
\]

Пример Ethernet-кадра:
\[
\text{Dst MAC} \;|\; \text{Src MAC} \;|\; \text{Type} \;|\; \text{Payload} \;|\; \text{FCS}.
\]

\subsection{Сетевой уровень}

Сетевой уровень отвечает за доставку пакетов между узлами, находящимися в разных сетях.

Основные функции:
\begin{itemize}
    \item логическая адресация;
    \item маршрутизация;
    \item выбор следующего перехода;
    \item фрагментация и сборка пакетов;
    \item диагностика ошибок.
\end{itemize}

Единица данных:
\[
\text{пакет}.
\]

Основной пример протокола сетевого уровня --- IP.

\subsection{Транспортный уровень}

Транспортный уровень обеспечивает взаимодействие процессов на конечных узлах.

Функции:
\begin{itemize}
    \item мультиплексирование приложений с помощью портов;
    \item сегментация данных;
    \item контроль доставки;
    \item управление потоком;
    \item управление перегрузкой;
    \item восстановление порядка данных.
\end{itemize}

Единица данных:
\[
\text{сегмент} \quad \text{или} \quad \text{датаграмма}.
\]

TCP обеспечивает надёжную байтовую потоковую передачу, UDP обеспечивает ненадёжную передачу датаграмм с минимальными накладными расходами.

\subsection{Сеансовый уровень}

Сеансовый уровень управляет логическими сеансами обмена:
\begin{itemize}
    \item установление и завершение сеанса;
    \item управление диалогом;
    \item синхронизация;
    \item восстановление после сбоев на уровне сеанса.
\end{itemize}

В современных сетях многие функции сеансового уровня реализуются библиотеками, прикладными протоколами или транспортным уровнем.

\subsection{Уровень представления}

Уровень представления отвечает за форму представления данных:
\begin{itemize}
    \item кодирование символов;
    \item сериализация;
    \item преобразование форматов;
    \item сжатие;
    \item шифрование и дешифрование.
\end{itemize}

Примеры задач:
\[
\text{UTF-8} \leftrightarrow \text{UTF-16},
\quad
\text{объект} \leftrightarrow \text{байтовый поток}.
\]

\subsection{Прикладной уровень}

Прикладной уровень предоставляет сетевые службы конечным приложениям.

Примеры:
\begin{itemize}
    \item HTTP/HTTPS --- передача гипертекста;
    \item DNS --- преобразование доменных имён в IP-адреса;
    \item SMTP, IMAP, POP3 --- электронная почта;
    \item SSH --- удалённый доступ;
    \item FTP, SFTP --- передача файлов.
\end{itemize}

\subsection{Сопоставление OSI и TCP/IP}

\begin{center}
\begin{tabular}{|p{0.35\textwidth}|p{0.35\textwidth}|}
\hline
\textbf{OSI} & \textbf{TCP/IP} \\
\hline
Прикладной, представления, сеансовый & Прикладной \\
\hline
Транспортный & Транспортный \\
\hline
Сетевой & Интернет-уровень \\
\hline
Канальный, физический & Уровень сетевого доступа \\
\hline
\end{tabular}
\end{center}

\section{Маршрутизация в сетях}

\subsection{Основные определения}

\begin{definition}
\textbf{Маршрутизация} --- процесс выбора пути передачи пакетов от источника к получателю через сеть.
\end{definition}

\begin{definition}
\textbf{Маршрут} --- информация о том, через какой следующий узел или интерфейс следует передавать пакеты к определённой сети назначения.
\end{definition}

\begin{definition}
\textbf{Маршрутизатор} --- сетевое устройство, принимающее решения о пересылке пакетов на основе адреса назначения и таблицы маршрутизации.
\end{definition}

Таблица маршрутизации обычно содержит:
\begin{itemize}
    \item сеть назначения;
    \item маску или длину префикса;
    \item следующий переход;
    \item выходной интерфейс;
    \item метрику;
    \item источник маршрута.
\end{itemize}

Пример:

\begin{center}
\begin{tabular}{|c|c|c|c|}
\hline
\textbf{Назначение} & \textbf{Префикс} & \textbf{Next hop} & \textbf{Метрика} \\
\hline
192.168.1.0 & /24 & directly connected & 0 \\
\hline
10.0.0.0 & /8 & 192.168.1.1 & 10 \\
\hline
0.0.0.0 & /0 & 192.168.1.254 & 100 \\
\hline
\end{tabular}
\end{center}

Маршрут \texttt{0.0.0.0/0} называется \textbf{маршрутом по умолчанию}.

\subsection{Принцип наибольшего совпадения префикса}

При выборе маршрута IP-маршрутизатор использует правило \textbf{longest prefix match}: выбирается маршрут с наиболее длинным префиксом, совпадающим с адресом назначения.

\begin{example}
Пусть в таблице есть маршруты:
\[
10.0.0.0/8,\quad 10.1.0.0/16,\quad 10.1.2.0/24.
\]
Для адреса назначения:
\[
10.1.2.15
\]
подходят все три маршрута, но будет выбран:
\[
10.1.2.0/24,
\]
поскольку он имеет самый длинный префикс.
\end{example}

\subsection{Статическая и динамическая маршрутизация}

\textbf{Статическая маршрутизация} задаётся администратором вручную.

Преимущества:
\begin{itemize}
    \item простота;
    \item предсказуемость;
    \item отсутствие служебного трафика протоколов маршрутизации.
\end{itemize}

Недостатки:
\begin{itemize}
    \item плохая масштабируемость;
    \item ручное изменение при отказах;
    \item риск ошибок конфигурации.
\end{itemize}

\textbf{Динамическая маршрутизация} использует протоколы, которые автоматически строят и обновляют таблицы маршрутизации.

Примеры:
\begin{itemize}
    \item RIP --- дистанционно-векторный протокол;
    \item OSPF --- протокол состояния каналов;
    \item IS-IS --- протокол состояния каналов;
    \item BGP --- междоменный протокол маршрутизации.
\end{itemize}

\subsection{Метрика маршрута}

\begin{definition}
\textbf{Метрика} --- числовая или составная характеристика маршрута, используемая для выбора лучшего пути.
\end{definition}

Возможные компоненты метрики:
\begin{itemize}
    \item число переходов;
    \item пропускная способность каналов;
    \item задержка;
    \item загрузка;
    \item надёжность;
    \item административная стоимость;
    \item политика маршрутизации.
\end{itemize}

\section{Алгоритм Дейкстры для поиска кратчайших путей}

\subsection{Постановка задачи}

Пусть задан взвешенный ориентированный или неориентированный граф:
\[
G = (V,E),
\]
где:
\begin{itemize}
    \item $V$ --- множество вершин;
    \item $E$ --- множество рёбер;
    \item каждому ребру $(u,v)$ сопоставлен неотрицательный вес $w(u,v) \geq 0$.
\end{itemize}

Требуется найти кратчайшие расстояния от начальной вершины $s$ до всех остальных вершин.

В контексте сетей:
\begin{itemize}
    \item вершины --- маршрутизаторы;
    \item рёбра --- каналы связи;
    \item веса --- стоимости каналов;
    \item кратчайший путь --- маршрут с минимальной суммарной метрикой.
\end{itemize}

\subsection{Идея алгоритма}

Алгоритм поддерживает:
\begin{itemize}
    \item множество $S$ вершин, для которых кратчайшее расстояние уже окончательно найдено;
    \item оценки расстояний $d[v]$ для всех вершин;
    \item предков $p[v]$ для восстановления пути.
\end{itemize}

На каждом шаге выбирается вершина вне $S$ с минимальной текущей оценкой расстояния. После этого выполняется релаксация всех исходящих из неё рёбер.

\begin{definition}
\textbf{Релаксация ребра} $(u,v)$ --- проверка, можно ли улучшить текущую оценку расстояния до $v$, проходя через $u$:
\[
\text{если } d[v] > d[u] + w(u,v),
\text{ то } d[v] := d[u] + w(u,v).
\]
\end{definition}

\subsection{Шаги алгоритма Дейкстры}

\begin{enumerate}
    \item Для всех вершин $v \in V$ положить:
    \[
    d[v] := \infty,\quad p[v] := \varnothing.
    \]
    \item Для начальной вершины:
    \[
    d[s] := 0.
    \]
    \item Положить множество обработанных вершин:
    \[
    S := \varnothing.
    \]
    \item Пока существуют необработанные вершины:
    \begin{enumerate}
        \item выбрать вершину $u \notin S$ с минимальным $d[u]$;
        \item добавить $u$ в $S$;
        \item для каждого ребра $(u,v)$ выполнить релаксацию:
        \[
        \text{если } d[v] > d[u] + w(u,v),
        \]
        то:
        \[
        d[v] := d[u] + w(u,v), \quad p[v] := u.
        \]
    \end{enumerate}
    \item Значения $d[v]$ являются длинами кратчайших путей от $s$.
\end{enumerate}

\subsection{Псевдокод}

\begin{lstlisting}
Dijkstra(G, s):
    for each vertex v in V:
        d[v] := infinity
        p[v] := NIL

    d[s] := 0
    Q := V

    while Q is not empty:
        u := vertex in Q with minimal d[u]
        remove u from Q

        for each edge (u, v):
            if d[v] > d[u] + w(u, v):
                d[v] := d[u] + w(u, v)
                p[v] := u

    return d, p
\end{lstlisting}

\subsection{Минимальный пример}

Рассмотрим граф:

\[
A \xleftrightarrow{1} B,\quad
A \xleftrightarrow{4} C,\quad
B \xleftrightarrow{2} C,\quad
B \xleftrightarrow{5} D,\quad
C \xleftrightarrow{1} D.
\]

\begin{center}
\begin{tikzpicture}[node distance=2cm, every node/.style={circle,draw}]
\node (A) {A};
\node[right of=A] (B) {B};
\node[below of=A] (C) {C};
\node[right of=C] (D) {D};

\draw (A) -- node[above,draw=none] {1} (B);
\draw (A) -- node[left,draw=none] {4} (C);
\draw (B) -- node[right,draw=none] {2} (C);
\draw (B) -- node[right,draw=none] {5} (D);
\draw (C) -- node[below,draw=none] {1} (D);
\end{tikzpicture}
\end{center}

Начальная вершина:
\[
s=A.
\]

Инициализация:
\[
d[A]=0,\quad d[B]=d[C]=d[D]=\infty.
\]

\begin{center}
\begin{tabular}{|c|c|c|c|c|c|}
\hline
\textbf{Шаг} & \textbf{Выбрана} & $d[A]$ & $d[B]$ & $d[C]$ & $d[D]$ \\
\hline
0 & --- & 0 & $\infty$ & $\infty$ & $\infty$ \\
\hline
1 & A & 0 & 1 & 4 & $\infty$ \\
\hline
2 & B & 0 & 1 & 3 & 6 \\
\hline
3 & C & 0 & 1 & 3 & 4 \\
\hline
4 & D & 0 & 1 & 3 & 4 \\
\hline
\end{tabular}
\end{center}

Кратчайшие пути:
\[
A \to B: 1,
\]
\[
A \to C: A \to B \to C,\quad 1+2=3,
\]
\[
A \to D: A \to B \to C \to D,\quad 1+2+1=4.
\]

\subsection{Теорема корректности алгоритма Дейкстры}

\begin{theorem}
Пусть $G=(V,E)$ --- граф с неотрицательными весами рёбер $w(e)\geq 0$, а $s \in V$ --- начальная вершина. Тогда алгоритм Дейкстры после завершения для каждой вершины $v \in V$ находит длину кратчайшего пути из $s$ в $v$.
\end{theorem}

\begin{proof}
Обозначим через $\delta(s,v)$ истинную длину кратчайшего пути из $s$ в $v$.

Докажем инвариант: после добавления вершины $u$ в множество $S$ значение $d[u]$ равно $\delta(s,u)$.

Доказательство проведём по индукции по числу вершин в $S$.

\textbf{База.}
Первой выбирается вершина $s$, так как $d[s]=0$, а для остальных вершин начальные оценки равны $\infty$. Очевидно,
\[
d[s]=0=\delta(s,s).
\]
Значит, инвариант верен для первой выбранной вершины.

\textbf{Индукционный переход.}
Пусть перед очередным шагом для всех вершин из $S$ уже выполнено:
\[
d[x]=\delta(s,x), \quad x \in S.
\]
Пусть алгоритм выбирает вершину $u \notin S$ с минимальным значением $d[u]$ среди всех необработанных вершин.

Предположим противное:
\[
d[u] > \delta(s,u).
\]
Рассмотрим кратчайший путь $P$ из $s$ в $u$. Поскольку $s \in S$, а $u \notin S$, на этом пути существует первое ребро $(x,y)$, такое что:
\[
x \in S,\quad y \notin S.
\]
Вершина $x$ уже обработана. По предположению индукции:
\[
d[x]=\delta(s,x).
\]
Когда алгоритм обрабатывал $x$, он выполнил релаксацию ребра $(x,y)$, поэтому:
\[
d[y] \leq d[x] + w(x,y).
\]
Так как $x$ лежит на кратчайшем пути к $u$, имеем:
\[
\delta(s,y)=\delta(s,x)+w(x,y).
\]
Следовательно:
\[
d[y] \leq \delta(s,y).
\]
Но оценки $d[\cdot]$ никогда не могут стать меньше истинных кратчайших расстояний, так как они всегда являются длинами некоторых реально найденных путей. Поэтому:
\[
d[y]=\delta(s,y).
\]

Поскольку веса рёбер неотрицательны, расстояние от $s$ до $y$ не превосходит расстояние от $s$ до $u$ вдоль пути $P$:
\[
\delta(s,y) \leq \delta(s,u).
\]
Следовательно:
\[
d[y] = \delta(s,y) \leq \delta(s,u) < d[u].
\]
Получили, что существует необработанная вершина $y$ с оценкой меньше $d[u]$, что противоречит выбору $u$ как вершины с минимальной оценкой.

Значит, предположение неверно, и:
\[
d[u]=\delta(s,u).
\]

По индукции инвариант верен для всех вершин, добавленных в $S$. После завершения алгоритма все достижимые вершины обработаны, и для них найдены истинные кратчайшие расстояния. Для недостижимых вершин значение остаётся $\infty$, что также корректно.
\end{proof}

\subsection{Сложность алгоритма Дейкстры}

Если выбирать минимум простым просмотром массива:
\[
O(|V|^2).
\]

Если использовать двоичную кучу:
\[
O((|V|+|E|)\log |V|).
\]

Если использовать фибоначчиеву кучу:
\[
O(|E| + |V|\log |V|).
\]

\section{Дистанционно-векторная маршрутизация}

\subsection{Идея}

В дистанционно-векторной маршрутизации каждый маршрутизатор хранит для каждой сети назначения:
\begin{itemize}
    \item расстояние до неё;
    \item следующий переход.
\end{itemize}

Маршрутизаторы периодически обмениваются со своими соседями векторами расстояний.

Основой является уравнение Беллмана--Форда:
\[
D_x(y) = \min_{v \in N(x)} \{ c(x,v) + D_v(y) \},
\]
где:
\begin{itemize}
    \item $D_x(y)$ --- оценка стоимости пути от $x$ до $y$;
    \item $N(x)$ --- множество соседей узла $x$;
    \item $c(x,v)$ --- стоимость канала от $x$ до соседа $v$;
    \item $D_v(y)$ --- оценка соседа $v$ до назначения $y$.
\end{itemize}

\subsection{Шаги дистанционно-векторного алгоритма}

Для каждого маршрутизатора $x$:

\begin{enumerate}
    \item Инициализировать расстояние до себя:
    \[
    D_x(x)=0.
    \]
    \item Для каждого соседа $v$ задать:
    \[
    D_x(v)=c(x,v).
    \]
    \item Для остальных узлов:
    \[
    D_x(y)=\infty.
    \]
    \item Периодически отправлять соседям свой вектор расстояний.
    \item При получении вектора от соседа $v$ пересчитать для каждого назначения $y$:
    \[
    D_x(y) := \min(D_x(y), c(x,v)+D_v(y)).
    \]
    \item Если значение изменилось, обновить таблицу маршрутизации и при необходимости разослать новый вектор соседям.
\end{enumerate}

\subsection{Минимальный пример}

Пусть есть три маршрутизатора:
\[
A \xleftrightarrow{1} B \xleftrightarrow{2} C,
\quad
A \xleftrightarrow{5} C.
\]

Начальная таблица у $A$:
\[
D_A(A)=0,\quad D_A(B)=1,\quad D_A(C)=5.
\]

После получения от $B$ информации:
\[
D_B(C)=2,
\]
маршрутизатор $A$ пересчитывает:
\[
D_A(C)=\min(5, D_A(B)+D_B(C))=\min(5,1+2)=3.
\]

Новый путь:
\[
A \to B \to C,
\]
стоимость:
\[
3.
\]

\subsection{Проблема счёта до бесконечности}

Дистанционно-векторные протоколы подвержены проблеме \textbf{count to infinity}. Она возникает, когда маршрутизаторы постепенно увеличивают метрику до недоступной сети, ошибочно полагая, что путь существует через соседа.

Методы борьбы:
\begin{itemize}
    \item \textbf{split horizon}: не сообщать соседу маршрут, полученный от него же;
    \item \textbf{poison reverse}: сообщать такой маршрут с бесконечной метрикой;
    \item \textbf{triggered updates}: отправлять обновления сразу после изменения;
    \item ограничение максимальной метрики.
\end{itemize}

\section{Маршрутизация по состоянию каналов}

\subsection{Идея}

В протоколах состояния каналов каждый маршрутизатор:
\begin{enumerate}
    \item определяет своих соседей;
    \item измеряет стоимость каналов до соседей;
    \item формирует объявление состояния каналов;
    \item рассылает его всем маршрутизаторам области;
    \item строит полную карту топологии;
    \item запускает алгоритм Дейкстры для вычисления кратчайших путей.
\end{enumerate}

Пример такого подхода --- OSPF.

\subsection{Сравнение подходов}

\begin{center}
\begin{tabular}{|p{0.25\textwidth}|p{0.32\textwidth}|p{0.32\textwidth}|}
\hline
\textbf{Критерий} & \textbf{Дистанционно-векторная} & \textbf{Состояние каналов} \\
\hline
Информация & Расстояния до сетей & Полная топология области \\
\hline
Обмен & С соседями & Со всеми маршрутизаторами области \\
\hline
Алгоритм & Беллман--Форд & Дейкстра \\
\hline
Сходимость & Обычно медленнее & Обычно быстрее \\
\hline
Память и CPU & Меньше & Больше \\
\hline
Пример & RIP & OSPF, IS-IS \\
\hline
\end{tabular}
\end{center}

\section{Междоменная маршрутизация и BGP}

\begin{definition}
\textbf{Автономная система} --- совокупность IP-сетей и маршрутизаторов под единым административным управлением, имеющая общую политику маршрутизации.
\end{definition}

\begin{definition}
\textbf{BGP} --- Border Gateway Protocol, основной протокол междоменной маршрутизации в Интернете.
\end{definition}

BGP относится к протоколам типа \textbf{path-vector}. Он передаёт не только информацию о достижимости префиксов, но и путь в виде последовательности автономных систем.

Пример атрибута:
\[
AS\_PATH = 64500\; 64496\; 64497.
\]

BGP выбирает маршруты не только по технической метрике, но и по политике:
\begin{itemize}
    \item локальное предпочтение;
    \item длина AS\_PATH;
    \item тип происхождения маршрута;
    \item MED;
    \item eBGP/iBGP;
    \item стоимость до next-hop.
\end{itemize}

\section{Управление потоками данных в сети}

\subsection{Понятие потока данных}

\begin{definition}
\textbf{Поток данных} --- последовательность пакетов или байтов, передаваемых между источником и получателем и объединённых общими признаками.
\end{definition}

В сетях управление потоками может пониматься в нескольких смыслах:
\begin{enumerate}
    \item \textbf{flow control} --- защита получателя от переполнения;
    \item \textbf{congestion control} --- защита сети от перегрузки;
    \item \textbf{traffic shaping} --- формирование трафика;
    \item \textbf{traffic policing} --- ограничение трафика;
    \item \textbf{QoS} --- дифференцированное обслуживание разных классов трафика.
\end{enumerate}

\subsection{Управление потоком на канальном уровне}

На канальном уровне управление потоком может использоваться, если отправитель способен передавать кадры быстрее, чем получатель их обрабатывает.

Простейший вариант --- протокол \textbf{Stop-and-Wait}.

\subsection{Алгоритм Stop-and-Wait}

Идея: отправитель передаёт один кадр и ждёт подтверждения.

Шаги:
\begin{enumerate}
    \item Отправитель посылает кадр с номером $0$ или $1$.
    \item Запускает таймер.
    \item Получатель принимает кадр.
    \item Если кадр корректен, получатель отправляет ACK.
    \item Если ACK пришёл до истечения таймера, отправитель передаёт следующий кадр.
    \item Если таймер истёк, отправитель повторяет передачу.
\end{enumerate}

Преимущества:
\begin{itemize}
    \item простота;
    \item малое потребление буферов.
\end{itemize}

Недостаток:
\begin{itemize}
    \item низкая эффективность при больших задержках.
\end{itemize}

Коэффициент использования канала приближённо:
\[
U = \frac{T_{tx}}{T_{tx} + 2T_{prop}},
\]
где:
\begin{itemize}
    \item $T_{tx}$ --- время передачи кадра;
    \item $T_{prop}$ --- задержка распространения в одну сторону.
\end{itemize}

\subsection{Скользящее окно}

\begin{definition}
\textbf{Скользящее окно} --- механизм, позволяющий отправителю передавать несколько кадров или сегментов без ожидания подтверждения каждого из них.
\end{definition}

Пусть размер окна равен $W$. Тогда отправитель может иметь до $W$ неподтверждённых кадров.

\[
\underbrace{1\;2\;3}_{\text{подтверждены}}
\quad
\underbrace{4\;5\;6\;7}_{\text{окно}}
\quad
8\;9\;10
\]

Когда приходит подтверждение для кадра $4$, окно сдвигается:

\[
1\;2\;3\;4
\quad
\underbrace{5\;6\;7\;8}_{\text{новое окно}}
\quad
9\;10.
\]

\subsection{Управление потоком в TCP}

TCP использует окно приёма:
\[
rwnd,
\]
которое получатель сообщает отправителю. Оно показывает, сколько байтов получатель готов принять.

Отправитель ограничивает количество неподтверждённых данных:
\[
\text{in\_flight} \leq rwnd.
\]

Если буфер получателя заполняется, он уменьшает advertised window. Если окно становится нулевым, отправитель временно прекращает передачу пользовательских данных и периодически выполняет проверку окна.

\subsection{Управление перегрузкой в TCP}

Управление потоком защищает получателя, а управление перегрузкой защищает сеть.

TCP вводит окно перегрузки:
\[
cwnd.
\]

Фактическое окно отправки:
\[
W = \min(rwnd, cwnd).
\]

Классические фазы TCP congestion control:
\begin{itemize}
    \item slow start;
    \item congestion avoidance;
    \item fast retransmit;
    \item fast recovery.
\end{itemize}

\subsection{Slow Start}

В фазе slow start окно перегрузки растёт экспоненциально: примерно удваивается за каждый RTT.

Шаги:
\begin{enumerate}
    \item Установить малое начальное значение $cwnd$.
    \item При получении ACK увеличить $cwnd$.
    \item Пока:
    \[
    cwnd < ssthresh,
    \]
    продолжать экспоненциальный рост.
    \item При достижении порога $ssthresh$ перейти к congestion avoidance.
\end{enumerate}

Пример роста:
\[
1,\;2,\;4,\;8,\;16\;\text{MSS}.
\]

\subsection{Congestion Avoidance}

В фазе congestion avoidance рост становится примерно линейным:
\[
cwnd \leftarrow cwnd + \frac{MSS^2}{cwnd}
\]
на каждый ACK, что даёт примерно:
\[
+1 MSS
\]
за один RTT.

\subsection{Реакция на потери}

Потеря пакета интерпретируется как признак перегрузки.

Типичные действия:
\begin{itemize}
    \item уменьшить порог:
    \[
    ssthresh := \frac{cwnd}{2};
    \]
    \item уменьшить $cwnd$;
    \item повторно передать потерянный сегмент.
\end{itemize}

Если потеря обнаружена по тайм-ауту, реакция обычно более жёсткая, чем при получении нескольких дублирующих ACK.

\section{Формирование и ограничение трафика}

\subsection{Leaky Bucket}

\begin{definition}
\textbf{Leaky Bucket} --- алгоритм сглаживания трафика, при котором пакеты выходят из очереди с постоянной скоростью.
\end{definition}

Аналогия: ведро с отверстием. Вода поступает неравномерно, но вытекает равномерно.

Шаги алгоритма:
\begin{enumerate}
    \item Создать очередь фиксированного размера.
    \item Входящие пакеты помещать в очередь.
    \item С постоянной скоростью извлекать пакеты из очереди и передавать дальше.
    \item Если очередь заполнена, новые пакеты отбрасывать или маркировать.
\end{enumerate}

Минимальный пример:

Пусть скорость выхода:
\[
r=1 \text{ пакет/с},
\]
размер очереди:
\[
B=3.
\]

Если за одну секунду пришло 5 пакетов, то:
\begin{itemize}
    \item 3 пакета помещаются в очередь;
    \item 2 пакета отбрасываются;
    \item далее пакеты выходят по одному в секунду.
\end{itemize}

\subsection{Token Bucket}

\begin{definition}
\textbf{Token Bucket} --- алгоритм ограничения трафика, допускающий кратковременные всплески, если накоплены токены.
\end{definition}

Правила:
\begin{itemize}
    \item токены добавляются в ведро со скоростью $r$;
    \item максимальное число токенов равно $B$;
    \item для передачи пакета нужно потратить токены;
    \item если токенов недостаточно, пакет задерживается или отбрасывается.
\end{itemize}

Шаги алгоритма:
\begin{enumerate}
    \item Инициализировать число токенов:
    \[
    tokens := B.
    \]
    \item Через каждый интервал времени добавить:
    \[
    r \cdot \Delta t
    \]
    токенов, но не более $B$.
    \item При поступлении пакета размера $L$ проверить:
    \[
    tokens \geq L.
    \]
    \item Если условие выполнено:
    \[
    tokens := tokens - L,
    \]
    пакет передаётся.
    \item Иначе пакет ожидает или отбрасывается.
\end{enumerate}

Минимальный пример:

Пусть:
\[
r=10 \text{ токенов/с},\quad B=30.
\]

Если сеть простаивала 3 секунды, накопится:
\[
\min(30, 10 \cdot 3)=30
\]
токенов.

После этого можно сразу передать всплеск объёмом до 30 условных единиц.

\subsection{Сравнение Leaky Bucket и Token Bucket}

\begin{center}
\begin{tabular}{|p{0.3\textwidth}|p{0.28\textwidth}|p{0.28\textwidth}|}
\hline
\textbf{Свойство} & \textbf{Leaky Bucket} & \textbf{Token Bucket} \\
\hline
Выходной поток & Почти постоянный & Может быть всплесками \\
\hline
Всплески & Подавляются & Допускаются при наличии токенов \\
\hline
Назначение & Сглаживание & Ограничение средней скорости \\
\hline
При переполнении & Очередь переполняется & Токены перестают накапливаться \\
\hline
\end{tabular}
\end{center}

\section{Очереди и качество обслуживания}

\subsection{QoS}

\begin{definition}
\textbf{QoS} --- Quality of Service, совокупность механизмов, позволяющих управлять качеством передачи трафика разных классов.
\end{definition}

QoS особенно важен для:
\begin{itemize}
    \item голосовой связи;
    \item видеоконференций;
    \item потокового видео;
    \item промышленных систем управления;
    \item критичных бизнес-приложений.
\end{itemize}

Основные параметры QoS:
\begin{itemize}
    \item пропускная способность;
    \item задержка;
    \item jitter;
    \item потери пакетов;
    \item приоритет обслуживания.
\end{itemize}

\subsection{FIFO}

\begin{definition}
\textbf{FIFO} --- First In, First Out, дисциплина очереди, при которой пакеты обслуживаются в порядке поступления.
\end{definition}

Преимущество:
\begin{itemize}
    \item простота.
\end{itemize}

Недостатки:
\begin{itemize}
    \item отсутствие приоритетов;
    \item чувствительность к перегрузкам;
    \item возможное увеличение задержек для чувствительного трафика.
\end{itemize}

\subsection{Priority Queuing}

\begin{definition}
\textbf{Приоритетная очередь} --- дисциплина обслуживания, при которой пакеты высокого приоритета передаются раньше пакетов низкого приоритета.
\end{definition}

Недостаток:
\begin{itemize}
    \item возможное голодание низкоприоритетного трафика.
\end{itemize}

\subsection{Weighted Fair Queuing}

\begin{definition}
\textbf{WFQ} --- Weighted Fair Queuing, дисциплина очередей, распределяющая пропускную способность между классами трафика пропорционально весам.
\end{definition}

Если есть классы с весами:
\[
w_1,w_2,\ldots,w_n,
\]
то доля полосы для класса $i$:
\[
\frac{w_i}{\sum_{j=1}^{n} w_j}.
\]

Пример:

Пусть веса трёх классов:
\[
w_1=1,\quad w_2=2,\quad w_3=3.
\]
Сумма:
\[
1+2+3=6.
\]
Тогда доли:
\[
\frac{1}{6},\quad \frac{2}{6},\quad \frac{3}{6}.
\]

Если канал имеет скорость 60 Мбит/с, то классы получат:
\[
10,\quad 20,\quad 30 \text{ Мбит/с}.
\]

\section{Маршрутизация и управление потоками: связь задач}

Маршрутизация и управление потоками решают разные, но связанные задачи.

\begin{center}
\begin{tabular}{|p{0.28\textwidth}|p{0.28\textwidth}|p{0.3\textwidth}|}
\hline
\textbf{Задача} & \textbf{Главный вопрос} & \textbf{Примеры механизмов} \\
\hline
Маршрутизация & Куда отправить пакет? & Таблицы маршрутизации, OSPF, BGP \\
\hline
Управление потоком & Сколько можно отправить получателю? & TCP receive window, sliding window \\
\hline
Управление перегрузкой & Сколько можно отправить, не перегружая сеть? & TCP cwnd, AIMD \\
\hline
QoS & Как обслуживать разные классы трафика? & Очереди, приоритеты, shaping \\
\hline
\end{tabular}
\end{center}

В реальной сети эти механизмы действуют совместно:
\begin{itemize}
    \item маршрутизаторы выбирают путь;
    \item очереди на интерфейсах определяют порядок передачи;
    \item TCP регулирует скорость отправителя;
    \item QoS задаёт приоритеты;
    \item системы мониторинга обнаруживают перегрузки и отказы.
\end{itemize}

\section{Краткие выводы}

\begin{enumerate}
    \item Операционные средства управления сетями включают функции ОС, сетевые службы, диагностические утилиты, мониторинг, журналирование и протоколы управления.
    \item Модель ISO/OSI разделяет сетевое взаимодействие на семь уровней: физический, канальный, сетевой, транспортный, сеансовый, представления и прикладной.
    \item Маршрутизация отвечает за выбор пути передачи пакетов через сеть.
    \item Статическая маршрутизация проста, но плохо масштабируется; динамическая маршрутизация использует специальные протоколы.
    \item Алгоритм Дейкстры применяется в протоколах состояния каналов и находит кратчайшие пути при неотрицательных весах.
    \item Дистанционно-векторная маршрутизация основана на уравнении Беллмана--Форда.
    \item BGP является основным протоколом междоменной маршрутизации и учитывает политики автономных систем.
    \item Управление потоками защищает получателя, управление перегрузкой защищает сеть, а QoS позволяет дифференцировать обслуживание трафика.
    \item Алгоритмы Leaky Bucket и Token Bucket применяются для сглаживания и ограничения трафика.
\end{enumerate}

\end{document}
